% ==================================================================
%  RELATED WORK + DISCUSSION (pp. 8--9)
% ==================================================================

\section{Related Work}
\label{sec:related}

\paragraph{Compound AI systems and routing.}
The shift from monolithic models to compound
systems~\citep{zaharia2024compound} has generated a growing literature on
LLM routing---selecting which model handles which query.
RouteLLM~\citep{ong2024routellm} and
FrugalGPT~\citep{chen2023frugalgpt} learn routers from preference data
or cascading strategies; xRouter and MasRouter train routing policies via
reinforcement learning or distillation.  These systems solve the routing \emph{execution}
problem---given a trained router, which model to call.  Our work addresses
the upstream \emph{design selection} problem: given a menu of possible
system configurations (model, depth, routing policy), which configuration
to deploy.  The four-term decomposition provides the analytical scaffold
that routing systems lack: it quantifies \emph{why} a routing change
improves performance, separating information gain from mismatch reduction.

\paragraph{Test-time compute scaling.}
The discovery that repeated sampling can compensate for model
size~\citep{brown2024monkeys,snell2025testtime} motivates the
crossover formula~\eqref{eq:crossover}.
\citet{wu2025inference} derive inference scaling laws relating compute
to accuracy; \citet{osca2024} optimize sample-compute allocation.
Our crossover formula makes the tradeoff explicit at the level of individual
task modes, rather than in aggregate.
The key distinction is that these works study \emph{how much}
test-time compute to use; our decomposition studies \emph{which system
configuration} to use, with test-time compute as one of several levers.

\paragraph{Algorithm selection and transfer.}
The algorithm selection problem~\citep{rice1976algorithm} asks which
solver to apply to a given problem instance.
SATzilla~\citep{xu2012satzilla} and AutoFolio~\citep{lindauer2015autofolio}
solve this with portfolio-based methods; SMAC~\citep{hutter2011smac} and
Hyperband~\citep{li2017hyperband} use sequential model-based or
bandit-based optimization.  These treat each solver as a black box.
Our separable evaluator exploits the packed-family structure to score
solvers analytically from shared pilot data, achieving $O(1)$-in-$k$
evaluation cost where portfolio methods require $O(k)$.

Transfer of algorithm configurations across problem
distributions connects to domain adaptation
theory~\citep{bendavid2010theory,mansour2009domain} and
transferability estimation~\citep{achille2019task2vec,you2021logme,
pandy2022transferability}.  The Ben-David theory bounds generalization
error under distribution shift via $\mathcal{H}$-divergence; our transfer bounds
use the more structured $(\alpha_{\mathrm{PT}}, \alpha_O)$ divergence
pair, which exploits the decomposition to separate competence shift from
routing shift.  This finer-grained analysis enables the three-case
decision rule, which has no analogue in generic domain adaptation.

Racing algorithms~\citep{maron1993hoeffding,maron1997racing} and
successive halving~\citep{jamieson2016successive} reduce evaluation cost
by eliminating poor candidates early.  They remain black-box methods:
each candidate must be evaluated individually.
Algorithm~\ref{alg:decomp-guided} is white-box in the sense that it
exploits the algebraic structure of the performance objective to avoid
per-candidate evaluation entirely in Phase~2.
Weitzman's Pandora's box rule~\citep{weitzman1979optimal} solves optimal
sequential search with inspection costs; our setting differs in that the
separable evaluator makes all inspections simultaneous.

\section{Discussion}
\label{sec:discussion}

\paragraph{Limitations.}
The packed-family assumption is the central modeling commitment.
It requires that task difficulty clusters into discrete modes with
geometric time distributions---a reasonable approximation for
verifiable tasks with independent attempts, but a simplification for
tasks with complex dependencies between retries (e.g., where each
attempt builds on the previous one).  Outside the packed family, the
four-term identity holds only approximately, with a bounded residual
that depends on departure from the geometric-trials
model~\citep{beneventano2026framework}.

The transfer thresholds are conservative by design.  The 73.8\% accuracy
on borderline cells reflects worst-case Pinsker and entropy-continuity
bounds; tighter bounds (e.g., Bernstein-type refinements or empirical
validation on diagnostic samples) would likely improve accuracy.  The
thresholds should be treated as guidelines, not hard boundaries.

The framework is most valuable when models have comparative advantages across modes and $k \gg m \cdot \Mn$; when one model dominates everywhere (as in EXP15), the decomposition correctly diagnoses $G \approx 0$ but offers no advantage over simple model selection.

The design space in our experiments is finite and low-dimensional
(model $\times$ depth).  Extending to continuous design spaces---joint
optimization over model mixtures, context budgets, and tool
configurations---is an important open direction.

\paragraph{Takeaways.}
Three things are clearer after this work.
First, the separation of cost, competence, information, and routing
is not merely a diagnostic convenience---it is the algebraic structure
that makes $O(1)$-in-$k$ design selection possible.
Second, the two dimensions of transfer (competence and orchestration)
are genuinely independent, which means that a single similarity score
between task domains is insufficient for transfer decisions.
Third, even when a framework's assumptions are violated (as in EXP15,
where routing is trivial), the decomposition correctly diagnoses the
violation and remains useful for scoping.

\paragraph{Open questions.}
Can the sample-complexity cascade be tightened beyond the Bernstein
refinement?  Empirical pilot sizes are roughly $29\times$ smaller than
the Hoeffding bound predicts; closing this gap would make the budget
prescriptions practical rather than merely conservative.
Can the framework be extended to tasks where attempts are not independent
(e.g., agentic loops with memory)?  The geometric-trials assumption is
the binding constraint.
Finally, validating the decomposition and transfer theory on production
agent deployments---where mode structure must be discovered rather than
assumed---is the most important next step.

\paragraph{Acknowledgements.}
This work was conducted using the PoggioAI/MSc research pipeline
\citep{PoggioAIMSc_2026}, an experimental framework for ML theory
research with humans on the loop.
